\section{Results}\label{sec:results}

\subsection{The filing predicts the disposal}\label{sec:gate1}

Table~\ref{tab:gate1} reports how well the disposal signal separates the subsidiaries that were later disposed of from those retained, with the area under the ROC curve as the measure (0.5 is chance, 1 is a perfect ordering). Across the \NGateOne{} clean-label divisions of 2001--2025 the two-model signal reaches \AUCPooledLLM{}. A logistic model on the financial characteristics stated in the filing, estimated out of sample with repeated cross-validation, reaches \AUCPooledFin{}; the difference of \DeltaPooled{} has a paired bootstrap interval of [\DeltaPooledLo, \DeltaPooledHi]. The signal is stronger in the recent period (\AUCNewLLM{} for 2017--2025 against \AUCNewFin{} for the financial model) than in the earlier ones. A financial model trained on the 2001--2016 divisions and scored on 2017--2025 reaches \AUCTemporalFin{} against the signal's \AUCTemporalLLM{} on the same test set, so the gap is not an artefact of fitting within period. The filing, in other words, says more about what the controller will do than the balance sheets of the two companies it creates.

Three filings show what the models read. Doosan's filing of June 2021 for the carve-out of its forklift business stated that, once the division was complete, the parent would sell its entire stake in the new company to an affiliate, Doosan Bobcat; the two-model signal is \ExDoosanScore{} and the subsidiary was \ExDoosanOutcome{}. CJ ENM's filing of July 2020 for the carve-out of its streaming service listed, among the division's purposes, ``attracting outside investment, strategic alliances, technological cooperation and the sale of shares where necessary''; the signal is \ExTvingScore{} and the subsidiary was \ExTvingOutcome{}. DB HiTek's filing of May 2023 stated that the parent would hold 100 percent of the new company and wrote into both companies' charters that any listing within five years of the division would require a special resolution of the parent's shareholders; the signal is \ExDBScore{} and the subsidiary was \ExDBOutcome{}. The sentences the models quote are the sentences a careful shareholder would have underlined.

The anonymized condition answers the recall objection. For the 2001--2016 divisions, where the models most often reported recognising the company, the two-model signal falls from \AUCOldIdent{} with the names to \AUCOldAnon{} without them; the paired bootstrap interval of the gap, [\GapOldLo, \GapOldHi], includes zero. For 2017--2026 it falls from \AUCNewIdent{} to \AUCNewAnon. The filings Model~A said it recognised are the ones it scores best, but it scores them equally well when the names are masked (\AUCClaudeFlaggedAnonOld{} against \AUCClaudeFlaggedIdentOld{}). Those filings are easier to read; their outcomes were not remembered. The return results below have been repeated with the anonymized signal in place of the identified one. The Spearman correlation with the one-year return is \SpearmanAnon{} ($p$ = \SpearmanAnonP), the within-year rank coefficient with the full controls \RankAnonCoef{} ($p$ = \RankAnonP), and the early-disposal AUC \AUCEarlyAnon{} (Supporting Information).

The signal is also sharper about the disposals that matter for the one-year return. Against the retained subsidiaries, it orders the subsidiaries disposed of within one fiscal year of the division with an AUC of \AUCEarlyPooled{} (\AUCEarlyNew{} in 2017--2026), and those disposed of two or more years later with \AUCLateVsNever. What the filing tells a reader is, above all, whether the controller already has a buyer, an investor or a listing in mind.

\begin{sidewaystable}[htbp]
\centering
\caption{Disposal prediction from the filing text: identified and anonymized conditions}
\label{tab:gate1}
{\footnotesize\setlength{\tabcolsep}{3pt}% generated by code/make_tables_cgir.py on 2026-10-28 — source: gate1_pooled_20261019.json (pooled, period_*, temporal_*); anonymized_check_20261025.json by_period
\begin{tabular}{lrrrrrrrc}
\toprule
 & & Disposed & Financial & \multicolumn{3}{c}{Identified filing text} & Anonymized & Gap 95\% CI \\
\cmidrule(lr){5-7}
Sample & $n$ & share (\%) & baseline & Model A & Model B & Two-model & two-model & (identified $-$ anonymized) \\
\midrule
2001--2008 & 161 & 77.0 & 0.400 & 0.625 & 0.631 & 0.639 & -- & -- \\
2009--2016 & 178 & 71.9 & 0.443 & 0.771 & 0.707 & 0.751 & -- & -- \\
2017--2025 & 240 & 62.1 & 0.577 & 0.813 & 0.804 & 0.813 & 0.803 & [$-$0.008, 0.025] \\
Pooled 2001--2025 & 579 & 69.3 & 0.587 & 0.771 & 0.736 & 0.761 & -- & -- \\
\quad 2001--2016, both conditions & 344 & -- & -- & 0.716 & 0.674 & 0.706 & 0.691 & [$-$0.006, 0.039] \\
\quad 2017--2026, both conditions & 232 & -- & -- & 0.812 & 0.803 & 0.810 & 0.803 & [$-$0.008, 0.025] \\
\midrule
Temporal split: train $\le$2016 ($n$=339), test $\ge$2017 ($n$=240) & & & 0.611 & & & 0.813 & & \\
\bottomrule
\end{tabular}
}
\tabnote{Area under the ROC curve for the clean-label divisions (disposed versus retained at 100 percent). The financial baseline is a logistic regression on characteristics stated in the filing (new company's assets, debt to assets and sales to assets; relative size; surviving company's debt to assets; listing application; market) estimated with repeated stratified five-fold cross-validation (50 repeats); year is excluded. Model~A and Model~B are two large language models reading the filing text as of its date; the two-model score is their average. The anonymized condition masks the parent's and the new company's names, group names and dates and withholds the filing date; the gap interval is a paired bootstrap over events. The temporal split trains the financial model on 2001--2016 and scores 2017--2025.}
\end{sidewaystable}

\subsection{The market's reaction and the year after}\label{sec:market}

Table~\ref{tab:returns_by_score} sorts the listed parents into quintiles of the disposal signal and reports the announcement return, the one-year return and the market-model cumulative return over the same year. Three facts are in the table. First, the announcement reaction does not fall with the signal; it rises. The slope of the three-day abnormal return on the signal is \AnnCarSlope{} (standard error \AnnCarSlopeSE{}, $p$ = \AnnCarSlopeP{}) with year fixed effects, so a filing that reads as a certain disposal is received \AnnCarSlopePP{} percentage points more favourably than one that reads as a certain retention. Second, the return over the following year falls with the signal. The Spearman correlation between the signal and the one-year return is \SpearmanScoreBhar{} ($p$ = \SpearmanScoreBharP) across \NScoreOneY{} parents, and \SpearmanWithinYear{} ($p$ = \SpearmanWithinYearP) when both are ranked within the announcement year. The difference between the top and bottom quintiles is \QQWithinYear{} percentage points in the within-year sort and \QQDiff{} points in the pooled sort, with wide intervals because the returns are heavy-tailed; at the median it is \QQMedianDiff{} points (\QFiveMedian{} against \QOneMedian{} percent; Mann--Whitney $p$ = \QQMWP). Third, the calendar-time tests point the same way without reaching conventional significance. A daily portfolio long the top tercile and short the bottom tercile earns \CTHiLo{} percent a year (Newey--West $t$ = \CTHiLoT), and its monthly five-factor alpha is \FFHiLo{} percent a month ($t$ = \FFHiLoT).

We read the three facts together. The signal is public information at the announcement; a reader who extracts it does better than the balance sheets at forecasting the disposal; the price reacts to it as to good news, not as to a transfer; and the parents it marks go on to underperform. The rank-based evidence for the last step is clear and the mean-based evidence is suggestive. That is what a right-skewed return distribution produces when the typical event loses and a few events gain a great deal. Across all \NScoreOneY{} parents the mean one-year market-adjusted return is \DriftMean{} percent (bootstrap interval [\DriftLo, \DriftHi]), while the median is \DriftMedian{} percent and \DriftNeg{} percent of parents trail their index; the five-factor alpha of the equal-weighted calendar-time portfolio is \FFAllAlpha{} percent a month ($t$ = \FFAllT).

\begin{sidewaystable}[htbp]
\centering
\caption{Announcement reaction and one-year returns by the filing's disposal signal}
\label{tab:returns_by_score}
{\footnotesize\setlength{\tabcolsep}{3pt}% generated by code/make_tables_cgir.py on 2026-10-28 — source: bhar_predicted_disposal_pooled_20261028.json: by_score_quintile, spearman_score_vs_return, regressions, calendar_time_daily_madj, calendar_time_monthly_ff5
\begin{tabular}{>{\raggedright\arraybackslash}p{5.6cm}rrrrrr}
\toprule
Disposal-score quintile & $n$ & \multicolumn{3}{c}{One-year market-adjusted BHAR [+2, +250] (\%)} & CAR$_{\text{mm}}$ [+2, +250] (\%) & Announcement CAR [$-$1, +1] (\%) \\
\cmidrule(lr){3-5}
 & & Mean & Winsorised mean & Median & Mean & Mean \\
\midrule
Q1 (lowest) & 119 & +1.5 & +0.2 & $-$3.0 & $-$9.0 & +0.49 \\
Q2 & 105 & +1.7 & $-$0.7 & $-$17.6 & $-$8.4 & +0.74 \\
Q3 & 112 & $-$12.1 & $-$13.2 & $-$19.2 & $-$16.8 & +0.11 \\
Q4 & 118 & $-$15.7 & $-$17.7 & $-$24.0 & $-$20.5 & +0.48 \\
Q5 (highest) & 105 & $-$2.9 & $-$10.0 & $-$23.0 & $-$7.1 & +1.26 \\
\midrule
Q5 $-$ Q1 & & $-$4.4 & & $-$20.0 & +1.9 & +0.78 \\
\quad bootstrap 95\% CI & & [$-$22.7, +17.3] & & & & [$-$1.30, +2.83] \\
\quad Mann--Whitney $p$ & & 0.008 & & & 0.702 & 0.593 \\
Within-year quintiles, Q5 $-$ Q1 & & $-$11.7 & & & & \\
Spearman $\rho$(score, outcome) [$p$] & 559 & $-$0.136 [0.001] & & & $-$0.009 [0.827] & 0.033 [0.438] \\
\quad within-year ranks & & $-$0.134 [0.002] & & & & \\
OLS slope on score, year FE [HC1 $p$] & 559 & $-$0.186 [0.155] & & & 0.002 [0.989] & 0.0372 [0.046] \\
\midrule
Calendar-time, top $-$ bottom tercile: annualised AR & & $-$12.2 & \multicolumn{2}{l}{NW $t$ = $-$1.68} & & \\
FF5 alpha, top $-$ bottom tercile (monthly) & & $-$0.52 & \multicolumn{2}{l}{NW $t$ = $-$0.96} & & \\
\bottomrule
\end{tabular}
}
\tabnote{Listed parents with a disposal signal and a full one-year window, 2001--2026, sorted into quintiles of the two-model signal. BHAR is the parent's buy-and-hold return minus its market index's over trading days $+2$ to $+250$; CAR$_{\text{mm}}$ sums market-model abnormal returns over the same days; the announcement CAR is the market-model abnormal return over days $-1$ to $+1$. Winsorised means clip at the second and ninety-eighth percentiles. Bootstrap intervals use 10,000 resamples. The OLS slope is the coefficient on the signal (0--1) with announcement-year fixed effects and HC1 standard errors. Calendar-time: equal-weighted daily portfolio of parents within their windows, top minus bottom tercile, annualised; the five-factor alpha is from a monthly regression on Korean market, size, value, profitability and investment factors with Newey--West standard errors.}
\end{sidewaystable}

Table~\ref{tab:returns_by_outcome} replaces the ex-ante signal with the realised outcome and its timing. The mean one-year return of parents whose subsidiary was disposed of differs from that of parents that retained it by \DispRetDiff{} percentage points (Mann--Whitney $p$ = \DispRetMWP; Welch $p$ = \DispRetWelchP), but the average conceals the timing. The \EarlyN{} parents whose subsidiary changed hands within one fiscal year of the division earned \EarlyMean{} percent on average and \EarlyMedian{} percent at the median, against \NeverMean{} and \NeverMedian{} percent for the \NeverN{} parents that retained; the difference in means is \EarlyNeverDiff{} points (Mann--Whitney $p$ = \EarlyNeverMWP; bootstrap interval [\EarlyNeverLo, \EarlyNeverHi]). Parents whose disposal came two or more years later earned \LateMean{} percent, indistinguishable from the retainers ($p$ = \LateNeverP). The loss is where \citeA{kang2025} located it, in the disposal, and it arrives when the disposal does.

\begin{table}[htbp]
\centering
\caption{One-year returns by realized outcome and its timing}
\label{tab:returns_by_outcome}
{\footnotesize\setlength{\tabcolsep}{3pt}% generated by code/make_tables_cgir.py on 2026-10-26 — source: bhar_predicted_disposal_pooled_20261025.json: by_realized_outcome, by_disposal_timing, calendar_time_*
\begin{tabular}{>{\raggedright\arraybackslash}p{5.3cm}rrrrr}
\toprule
Realized outcome & $n$ & Mean (\%) & Winsorised mean (\%) & Median (\%) & Share negative (\%) \\
\midrule
Disposed (any time) & 330 & $-$8.9 & $-$10.7 & $-$22.5 & 70 \\
Retained at 100\% & 141 & +1.4 & $-$2.6 & $-$8.7 & 60 \\
\quad disposed $-$ retained & & $-$10.3 & & $-$13.8 & \\
\quad bootstrap 95\% CI; Welch $p$; Mann--Whitney $p$ & & [$-$27.9, +4.9] & 0.219 & 0.010 & \\
\midrule
Early disposal (first change within one fiscal year) & 195 & $-$14.2 & $-$16.0 & $-$26.6 & 72 \\
Late disposal (two or more fiscal years) & 70 & $-$3.8 & $-$5.4 & $-$19.9 & 69 \\
Disposed, timing not recorded & 65 & +1.3 & $-$1.5 & $-$14.5 & 65 \\
Never (retained at 100\%) & 141 & +1.4 & $-$2.6 & $-$8.7 & 60 \\
\quad early $-$ never & & $-$15.6 & & $-$17.9 & \\
\quad bootstrap 95\% CI; Welch $p$; Mann--Whitney $p$ & & [$-$34.2, +1.3] & 0.082 & 0.001 & \\
\quad late $-$ never & & $-$5.2 & & & \\
\quad Welch $p$; Mann--Whitney $p$ & & & 0.626 & 0.236 & \\
\midrule
Calendar-time, disposed $-$ retained: annualised AR (NW $t$) & & $-$2.9 & ($-$0.32) & & \\
FF5 alpha, disposed $-$ retained, monthly (NW $t$) & & +0.19 & (0.25) & & \\
\bottomrule
\end{tabular}
}
\tabnote{Listed parents with a clean label and a full one-year window. Early disposal: the parent's annual report first shows the new company sold, listed or diluted in the fiscal year of the division or the next; late: two or more fiscal years later; timing not recorded: disposed subsidiaries whose first change is known only from the equity-investment schedule. Returns as in Table~\ref{tab:returns_by_score}. Calendar-time and five-factor rows compare the disposed and retained portfolios.}
\end{table}

Table~\ref{tab:controls} asks whether the association between the signal and the year's return is the parent's condition in disguise: a troubled parent might both write a filing that reads as a disposal and go on to underperform. The dependent variable in the first row is the parent's percentile rank of the one-year return within its announcement year, which is immune to the skew. The coefficient on the signal is \RankSpecOneCoef{} with year fixed effects alone, \RankSpecTwoCoef{} when size, book-to-market, the six-month pre-announcement return and the market are added, and \RankSpecThreeCoef{} when the new company's relative size, the surviving company's debt ratio and the stated listing plan are added as well (HC1 $p$ = \RankSpecThreeP; $p$ \RankSpecThreeClP{} clustered by year); across the three pooled specifications it lies between \RankScoreCoefMin{} and \RankScoreCoefMax{}. A filing that reads as a certain disposal rather than a certain retention thus places the parent \RankSpecThreePctl{} percentiles lower in its year's return ranking. The coefficient is \RankSpecThreeBCoef{} in 2001--2016 alone ($p$ = \RankSpecThreeBP), \RankSpecThreeCCoef{} in 2017--2026 alone ($p$ = \RankSpecThreeCP) and \RankSpecFourCoef{} in 2017--2026 with the parent's own return on assets, debt ratio and size added ($p$ = \RankSpecFourP). The coefficients on the mean return in the second row are negative in five of the six specifications and imprecise in all. The slope of the announcement return on the signal in the fourth row stays positive in every specification; it is significant in the pooled sample and in 2001--2016, and not in 2017--2026 alone. The parent's condition does not explain the pattern.

\begin{sidewaystable}[htbp]
\centering
\caption{Returns on the disposal signal with controls}
\label{tab:controls}
{\footnotesize\setlength{\tabcolsep}{3pt}% generated by code/make_tables_cgir.py on 2026-10-28 — source: bhar_controls_20261025.json specs[bhar1y_rank|bhar1y|ann_car][1..4]
\begin{tabular}{>{\raggedright\arraybackslash}p{5.4cm}cccccc}
\toprule
Dependent variable & (1) & (2) & (3) & (4) & (5) & (6) \\
\midrule
Within-year percentile rank of one-year BHAR & $-$0.191*** (0.067) & $-$0.192*** (0.066) & $-$0.195*** (0.068) & $-$0.229** (0.098) & $-$0.168* (0.101) & $-$0.207** (0.104) \\
\quad $n$ & 559 & 559 & 559 & 350 & 209 & 209 \\
One-year market-adjusted BHAR & $-$0.187 (0.132) & $-$0.172 (0.133) & $-$0.158 (0.147) & $-$0.334 (0.205) & 0.010 (0.223) & $-$0.123 (0.253) \\
Announcement CAR [$-$1, +1] & 0.038** (0.019) & 0.041** (0.019) & 0.039** (0.019) & 0.060* (0.033) & 0.021 (0.022) & 0.027 (0.023) \\
\quad $n$ & 560 & 560 & 560 & 349 & 211 & 211 \\
\midrule
Year fixed effects & Yes & Yes & Yes & Yes & Yes & Yes \\
Size, book-to-market, pre-event return, market & & Yes & Yes & Yes & Yes & Yes \\
Filing: relative size, surviving leverage, listing plan & & & Yes & Yes & Yes & Yes \\
Parent fundamentals (2017+ only) & & & & & & Yes \\
Sample & Pooled & Pooled & Pooled & 2001--2016 & 2017--2026 & 2017--2026 \\
\bottomrule
\end{tabular}
}
\tabnote{Each cell is the OLS coefficient on the disposal signal (0--1) with the HC1 standard error in parentheses; $^{***}$, $^{**}$ and $^{*}$ mark $p<0.01$, $p<0.05$ and $p<0.10$. Columns (1)--(3) use the pooled sample, (4) 2001--2016, (5)--(6) 2017--2026. Controls: log market capitalisation at the month-end before the announcement and log book-to-market (Compustat Global); market-adjusted return over trading days $-130$ to $-2$; KOSDAQ indicator; the new company's assets relative to the surviving company's; the surviving company's debt to assets; whether the filing states a listing application; in (6) the parent's return on assets, debt to assets and log assets in the fiscal year before the division. Missing controls are set to the sample median with an indicator. All regressions include announcement-year fixed effects.}
\end{sidewaystable}

Table~\ref{tab:realized_controls} closes the loop between the signal and the realised disposal. The coefficient on early disposal in the one-year return is \JointEarlyCoef{} with the same controls ($p$ = \JointEarlyP), and when early disposal and the signal enter together the signal's coefficient is \JointScoreCoef{} ($p$ = \JointScoreP): once we know that the subsidiary changed hands within the year, what the filing said adds nothing to the return. The signal matters for the return because it predicts the early disposal, and the early disposal is where the minority's loss is realised.

\begin{table}[htbp]
\centering
\caption{Returns on realized early disposal}
\label{tab:realized_controls}
{\footnotesize\setlength{\tabcolsep}{3pt}% generated by code/make_tables_cgir.py on 2026-10-26 — source: bhar_controls_20261025.json specs.bhar1y_realized
\begin{tabular}{lcccc}
\toprule
 & (1) BHAR & (2) BHAR & (3) BHAR & (4) Rank \\
\midrule
Early disposal (within one fiscal year) & $-$0.143** & $-$0.132* & $-$0.151* & $-$0.079** \\
\quad HC1 $p$ & 0.021 & 0.057 & 0.083 & 0.013 \\
Disposal score & & & 0.135 & \\
\quad HC1 $p$ & & & 0.549 & \\
Controls, year FE & FE only & Yes & Yes & Yes \\
$n$ & 406 & 406 & 405 & 406 \\
\bottomrule
\end{tabular}
}
\tabnote{Clean-label listed parents with a full one-year window and a recorded disposal timing. Early disposal equals one when the first ownership change is in the fiscal year of the division or the next, zero for retained subsidiaries and late disposals. Controls and fixed effects as in Table~\ref{tab:controls}, columns (3); column (4) uses the within-year percentile rank of the one-year return.}
\end{table}

\subsection{The controlling owner's stake}\label{sec:block}

Tables~\ref{tab:block_terciles} and~\ref{tab:block_regressions} bring in the controlling block. Sorted into terciles of the block held before the announcement, the parents in the lowest tercile (blocks up to \LowBlockMax{} percent of common shares) filed divisions that read more strongly as disposals. They disposed of the subsidiary more often (\LowBlockDisp{} percent against \HighBlockDisp{} percent in the highest tercile, blocks of \HighBlockMin{} percent and above) and disposed of it early more often (\LowBlockEarly{} against \HighBlockEarly{} percent). And they earned less over the following year: a median one-year return of \LowBlockMedian{} percent against \HighBlockMedian{} percent (Mann--Whitney $p$ \LowHighMWP), a mean of \LowBlockMean{} against \HighBlockMean{} percent. The regressions put numbers on the chain. The disposal signal falls with the block (\ScoreOnBlock{} per unit of the block, HC1 $p$ = \ScoreOnBlockP, $p$ = \ScoreOnBlockClP{} clustered by year): controllers with smaller stakes write filings that a reader scores as more likely to end in a disposal. The realised outcomes move the same way with less precision: the coefficient of the disposal indicator on the block is \BlockOnDisposed{} (standard error \BlockOnDisposedSE) and that of the early-disposal indicator \BlockOnEarlyAlone{} (\BlockOnEarlyAloneSE) with year effects alone. Once the signal is added, the block's coefficient on early disposal is \BlockOnEarlyWithScore{} (\BlockOnEarlyWithScoreSE) while the signal's is \ScoreOnEarlyWithBlock{}. What the block tells us about the disposal is already in the filing. The parent's rank in its year's return rises with the block (\RankOnBlock{} per unit, $p$ = \RankOnBlockP). The interaction of early disposal with a below-median block shows that the two sources of loss overlap rather than add. In parents with a large block the coefficient of early disposal is \EarlyCoefBhar{} on the one-year return and \EarlyCoefRank{} on the within-year rank. That of a small block is \LowBlockCoefBhar{} and \LowBlockCoefRank{} regardless of disposal. The interaction terms, \InterCoefBhar{} ($p$ = \InterPBhar) and \InterCoefRank{} ($p$ = \InterPRank), offset most of the early-disposal effect in small-block parents, for whom the net coefficient of early disposal on the rank is \EarlyInLowBlockRank. Small-block parents trail whether or not the disposal has yet occurred; large-block parents trail when it has. The announcement reaction does not vary with the block (\BlockOnCar{} per unit, standard error \BlockOnCarSE).

This is the pattern the principal--principal frame predicts and the one \citeA{kang2025} found with a different sample and a different ownership measure: the minority shareholders who bear the loss are the ones whose controller had the least to lose. It is not a causal estimate. The block is measured before the announcement and cannot be caused by the division, but controllers with small stakes may run different businesses, and we have no instrument for the stake.

\begin{table}[htbp]
\centering
\caption{Carve-out outcomes and returns by the controlling block before the announcement}
\label{tab:block_terciles}
{\footnotesize\setlength{\tabcolsep}{3pt}% generated by code/make_tables_cgir.py on 2026-10-26 — source: governance_cross_section_20261025.json T5_block_terciles, T5_low_vs_high
\begin{tabular}{lrrr}
\toprule
 & Low block & Middle & High block \\
\midrule
Block range (\% of common shares) & 0--27 & 27--43 & 43--100 \\
$n$ & 184 & 183 & 184 \\
Disposal score (mean) & 0.550 & 0.516 & 0.506 \\
Disposal rate, clean labels (\%) & 74.0 & 70.6 & 61.5 \\
Early-disposal rate (\%) & 52.3 & 43.3 & 42.7 \\
Announcement CAR, mean (\%) & +1.04 & +0.80 & +0.95 \\
One-year BHAR, mean (\%) & $-$13.7 & +0.5 & $-$0.5 \\
One-year BHAR, median (\%) & $-$26.9 & $-$12.3 & $-$9.1 \\
\midrule
Low $-$ high: BHAR mean difference (\%); bootstrap 95\% CI & \multicolumn{3}{l}{$-$13.1; [$-$28.1, +1.4]} \\
\quad Welch $p$; Mann--Whitney $p$ & \multicolumn{3}{l}{0.081; $<$0.001} \\
\bottomrule
\end{tabular}
}
\tabnote{Listed parents with a pre-announcement ownership record (\NOwnership{} events, 2001--2026), sorted into terciles of the controlling block: common shares held by the largest shareholder and the related parties listed with it, excluding treasury shares, at the fiscal year-end before the announcement year. Disposal rates use clean labels; the early-disposal rate uses the subset with a recorded timing. Returns as in Table~\ref{tab:returns_by_score}.}
\end{table}

\begin{sidewaystable}[htbp]
\centering
\caption{The controlling block and the carve-out chain: regressions}
\label{tab:block_regressions}
{\footnotesize\setlength{\tabcolsep}{3pt}% generated by code/make_tables_cgir.py on 2026-10-28 — source: governance_cross_section_20261025.json T1..T4 (controls: size, B/M, pre-event return, market, filing variables; year FE; HC1 SE)
\begin{tabular}{lcccccc}
\toprule
Dependent variable & Block & Disposal score & Low block & Early disposal & Early $\times$ low block & $n$ \\
\midrule
Disposal score & $-$0.113** (0.045) & & & & & 527 \\
Disposed (any time) & $-$0.092 (0.129) & & & & & 446 \\
Early disposal & 0.113 (0.133) & 1.235*** (0.095) & & & & 382 \\
Announcement CAR [$-$1, +1] & 0.030 (0.023) & & & & & 537 \\
One-year BHAR & 0.120 (0.177) & & & & & 537 \\
Within-year rank of BHAR & 0.142* (0.076) & & & & & 537 \\
One-year BHAR & & & $-$0.231** (0.111) & $-$0.243** (0.096) & 0.223 (0.141) & 377 \\
Within-year rank of BHAR & & & $-$0.111*** (0.041) & $-$0.136*** (0.045) & 0.114* (0.063) & 377 \\
\bottomrule
\end{tabular}
}
\tabnote{OLS coefficients with HC1 standard errors in parentheses; $^{***}$, $^{**}$ and $^{*}$ mark $p<0.01$, $p<0.05$ and $p<0.10$. The block is a share (0--1); low block equals one below the sample median. Each row is one regression on the sample of Table~\ref{tab:block_terciles} with the dependent variable named in the first column, announcement-year fixed effects and the controls of Table~\ref{tab:controls}, column (3). Disposed and early disposal are indicators on clean labels.}
\end{sidewaystable}

\subsection{The 2022 reform}\label{sec:reform}

Table~\ref{tab:reform} compares the filings received before the revised disclosure form, from 2017 to September 2022, with the \NPostFilings{} received from October 2022 onward. The form did what forms do. A shareholder-protection statement appears in \ProtectPre{} percent of the earlier filings (one filing, dated 23 September 2022, between the package and the form) and in \ProtectPost{} percent of the later ones; a statement about the new company's listing plans appears in \ListingPre{} and \ListingPost{} percent. Captive spin-offs became rarer, from \FilingsPerYearPre{} filings a year before the cutoff to \FilingsPerYearPost{} after it. The announcement reaction changed sign and size, from \CarPre{} percent to \CarPost{} percent (Mann--Whitney $p$ \CarPrePostP), so that a captive spin-off announced under the new form is received as good news. What the filing signals about the disposal did not change: the mean disposal signal is \ScorePre{} before and \ScorePost{} after ($p$ = \ScorePrePostP). Nor, over the horizon we can see, did what the controllers did: among divisions old enough to have two fiscal years of annual reports, \EarlyPre{} percent disposed within a year before the reform and \EarlyPost{} percent after (Fisher $p$ = \EarlyPrePostP). The share eventually disposed of is lower after the reform, but that comparison is not fair to the later divisions, whose disposals have had less time to appear, and we do not rest on it.

Two cautions belong here. The post-reform sample is small (\NPost{} announcement returns, and fewer mature outcomes), and there is no group of divisions the reform did not touch. These are before-and-after differences, not estimates of the reform's effect. And the appraisal right took effect only in late December 2022, so the first post-reform quarter had the disclosure rule alone. With those limits, the reading is that the reform changed the filings and the market's reception of them within weeks, and that over the horizon we observe it did not change what a reader of the filing would forecast or how often the forecast came true.

\begin{sidewaystable}[htbp]
\centering
\caption{Filings, reactions and outcomes before and after the 2022 reform}
\label{tab:reform}
{\footnotesize\setlength{\tabcolsep}{3pt}% generated by code/make_tables_cgir.py on 2026-10-26 — source: regulation_2022_20261025.json pre_post (cutoff 2022-10-01)
\begin{tabular}{>{\raggedright\arraybackslash}p{6.0cm}rrr>{\raggedright\arraybackslash}p{2.6cm}}
\toprule
 & $n$ pre / post & Before (2017--Sep 2022) & After (Oct 2022--2026) & Test \\
\midrule
Filing states shareholder-protection measures (\%) & 208 / 44 & 0.5 & 95.5 & Fisher $<$0.001 \\
Filing mentions a listing plan for the new company (\%) & 208 / 44 & 9.1 & 77.3 & Fisher $<$0.001 \\
Stated purpose includes sale or disposal (\%) & 208 / 44 & 13.0 & 13.6 & Fisher 1.000 \\
Stated purpose includes fundraising or investor attraction (\%) & 208 / 44 & 15.9 & 27.3 & Fisher 0.084 \\
Disposal score (mean) & 208 / 44 & 0.528 & 0.510 & Welch 0.616; MW 0.442 \\
Announcement CAR [$-$1, +1] (\%) & 183 / 49 & +0.06 & +5.29 & Welch $<$0.001; MW $<$0.001 \\
Early disposal within one fiscal year, division year $\le$2023 (\%) & 117 / 19 & 40.2 & 36.8 & Fisher 1.000 \\
Disposed by 2025, division year $\le$2023 (\%) --- post period immature & 188 / 20 & 68.1 & 40.0 & Fisher 0.024 \\
Captive spin-off filings per year & & 39.0 & 16.4 & \\
\bottomrule
\end{tabular}
}
\tabnote{Captive spin-off filings by filers covered by the regulator's structured record, 2017 to 2026, split at 1 October 2022, the month in which the revised disclosure form took effect; the first filing after the cutoff to state protection measures is dated 17 October 2022. Filing content is read by the two models; a statement counts as present when the majority of the model judgments of the filing flag it. The disposal score is the two-model signal. Announcement CAR as in Table~\ref{tab:returns_by_score}. Outcome rates use divisions with at least two fiscal years of annual reports (division year 2023 or earlier). Fisher's exact test for shares; Welch and Mann--Whitney tests for means.}
\end{sidewaystable}
